\begin{theorem}[Chebyshev’s sum inequality]
Let $n\ge 1$ and let
\[
a_{1}\le a_{2}\le\cdots\le a_{n},\qquad
b_{1}\le b_{2}\le\cdots\le b_{n}
\]
be real numbers.  Then
\[
n\sum_{i=1}^{n}a_i b_i\;\ge\;
\Bigl(\sum_{i=1}^{n}a_i\Bigr)\Bigl(\sum_{i=1}^{n}b_i\Bigr).
\tag{1}
\]
\end{theorem}

\begin{proof}
\textbf{Step 1 – Trivial cases.}  
If $n=1$ the statement reduces to $a_{1}b_{1}\ge a_{1}b_{1}$, which holds with equality.  
If all $a_i$ are equal (i.e.\ $a_{1}=\dots =a_{n}$) then both sides of \eqref{1} are equal to $(\sum a_i)(\sum b_i)/n$; the same holds when all $b_i$ are equal.  These are the degenerate cases that will later appear as the equality conditions.

Henceforth assume $n\ge2$ and that the two sequences are not both constant.

\medskip
\textbf{Step 2 – Centering the sequences.}  
Define the arithmetic means
\[
\bar a:=\frac1n\sum_{i=1}^{n}a_i,\qquad
\bar b:=\frac1n\sum_{i=1}^{n}b_i .
\]

Write $a_i=\bar a+\alpha_i$ and $b_i=\bar b+\beta_i$ where
\[
\alpha_i:=a_i-\bar a,\qquad \beta_i:=b_i-\bar b,
\quad\text{so that}\quad
\sum_{i=1}^{n}\alpha_i=
\sum_{i=1}^{n}\beta_i=0 .
\]

Then
\begin{align*}
n\sum_{i=1}^{n}a_i b_i
&=n\sum_{i=1}^{n}(\bar a+\alpha_i)(\bar b+\beta_i)   \\
&=n\Bigl( n\bar a\bar b
        +\bar a\sum_{i=1}^{n}\beta_i
        +\bar b\sum_{i=1}^{n}\alpha_i
        +\sum_{i=1}^{n}\alpha_i\beta_i\Bigr)                \\
&=n^{2}\bar a\bar b + n\sum_{i=1}^{n}\alpha_i\beta_i .
\end{align*}
Since $\sum a_i = n\bar a$ and $\sum b_i = n\bar b$,
\[
\Bigl(\sum_{i=1}^{n}a_i\Bigr)\Bigl(\sum_{i=1}^{n}b_i\Bigr)
= n^{2}\bar a\bar b .
\]

Subtracting gives the identity
\[
n\sum_{i=1}^{n}a_i b_i-
\Bigl(\sum_{i=1}^{n}a_i\Bigr)\Bigl(\sum_{i=1}^{n}b_i\Bigr)
= n\sum_{i=1}^{n}\alpha_i\beta_i . \tag{2}
\]

Thus \eqref{1} is equivalent to
\[
\sum_{i=1}^{n}\alpha_i\beta_i\;\ge\;0 . \tag{2'}
\]

\medskip
\textbf{Step 3 – Covariance as a double sum.}  
We claim that
\[
\sum_{i=1}^{n}\alpha_i\beta_i
=\frac{1}{2n}\sum_{i=1}^{n}\sum_{j=1}^{n}
   (a_i-a_j)(b_i-b_j). \tag{3}
\]

\emph{Proof of \eqref{3}.}  
Expanding the right‑hand side,
\begin{align*}
\frac{1}{2n}\sum_{i,j}(a_i-a_j)(b_i-b_j)
&=\frac{1}{2n}\sum_{i,j}
   \bigl(a_i b_i - a_i b_j - a_j b_i + a_j b_j\bigr)  \\
&=\frac{1}{2n}
   \Bigl(
        n\sum_{i}a_i b_i
        -\sum_{i}a_i\sum_{j}b_j
        -\sum_{j}a_j\sum_{i}b_i
        + n\sum_{j}a_j b_j
   \Bigr)                                              \\
&=\frac{1}{2n}\bigl(2n\sum_{i}a_i b_i-2\sum_{i}a_i\sum_{j}b_j\bigr)\\
&=\sum_{i}a_i b_i-\frac{1}{n}\Bigl(\sum_{i}a_i\Bigr)\Bigl(\sum_{j}b_j\Bigr).
\end{align*}
Since $\bar a=\frac1n\sum a_i$ and $\bar b=\frac1n\sum b_i$, the last line equals
\[
\sum_{i}(a_i-\bar a)(b_i-\bar b)=\sum_{i}\alpha_i\beta_i .
\]
Thus \eqref{3} holds. \hfill$\square$

Combining \eqref{2} and \eqref{3} we obtain
\[
n\sum_{i=1}^{n}a_i b_i-
\Bigl(\sum_{i=1}^{n}a_i\Bigr)\Bigl(\sum_{i=1}^{n}b_i\Bigr)
= \frac{1}{2}\sum_{i=1}^{n}\sum_{j=1}^{n}
   (a_i-a_j)(b_i-b_j). \tag{4}
\]

\medskip
\textbf{Step 4 – Sign of each summand.}  
Because the two sequences are sorted in the same non‑decreasing order, for any indices $i,j$ we have
\[
i<j\;\Longrightarrow\; a_i\le a_j\ \text{and}\ b_i\le b_j,
\]
and similarly when $i>j$.  Consequently
\[
(a_i-a_j)(b_i-b_j)\ge 0\qquad\text{for all }i,j. \tag{5}
\]

Each term in the double sum of \eqref{4} is non‑negative, hence the whole double sum is non‑negative.  From \eqref{4} we deduce
\[
n\sum_{i=1}^{n}a_i b_i-
\Bigl(\sum_{i=1}^{n}a_i\Bigr)\Bigl(\sum_{i=1}^{n}b_i\Bigr)\ge0,
\]
which is precisely inequality \eqref{1}.

\medskip
\textbf{Step 5 – Equality cases.}  
Equality in \eqref{1} occurs exactly when the double sum in \eqref{4} is zero.  By \eqref{5} this happens iff every factor $(a_i-a_j)(b_i-b_j)$ is zero.  Hence for every pair $i,j$ we must have either $a_i=a_j$ or $b_i=b_j$.  The only possibilities compatible with the monotonicity of the two sequences are:

\begin{itemize}
\item all $a_i$ are equal (the sequence $(a_i)$ is constant), or
\item all $b_i$ are equal (the sequence $(b_i)$ is constant).
\end{itemize}
The case $n=1$ is contained in both alternatives.  In all other situations the double sum is strictly positive, and the inequality is strict.

\medskip
\textbf{Conclusion.}  
For any integer $n\ge1$ and any two real sequences that are similarly ordered, the Chebyshev sum inequality \eqref{1} holds.  Equality holds iff at least one of the sequences is constant (including the trivial $n=1$ case). \qed
\end{proof}